%\subsection{聖靈}

\subsection{那日}

\subsubsection{以賽亞書 2：2-22}
2\underline{ 末後的日子}，耶和華殿的山必堅立超乎諸山，高舉過於萬嶺，萬民都要流歸這山。 3 必有許多國的民前往，說：「來吧！我們登耶和華的山，奔雅各神的殿。主必將他的道教訓我們，我們也要行他的路，因為訓誨必出於錫安，耶和華的言語必出於耶路撒冷。」 4 他必在列國中施行審判，為許多國民斷定是非。他們要將刀打成犁頭，把槍打成鐮刀。這國不舉刀攻擊那國，他們也不再學習戰事。5 雅各家啊，來吧，我們在耶和華的光明中行走！ 6 耶和華，你離棄了你百姓雅各家，是因他們充滿了東方的風俗，做觀兆的，像非利士人一樣，並與外邦人擊掌。 7 他們的國滿了金銀，財寶也無窮；他們的地滿了馬匹，車輛也無數。 8 他們的地滿了偶像，他們跪拜自己手所造的，就是自己指頭所做的。 9 卑賤人屈膝，尊貴人下跪，所以不可饒恕他們。 10 你當進入巖穴，藏在土中，躲避耶和華的驚嚇和他威嚴的榮光。 11 到\underline{ 那日}，眼目高傲的必降為卑，性情狂傲的都必屈膝，唯獨耶和華被尊崇。12 必有萬軍耶和華\underline{ 降罰的一個日子}，要臨到驕傲狂妄的，一切自高的都必降為卑； 13 又臨到黎巴嫩高大的香柏樹和巴珊的橡樹； 14 又臨到一切高山的峻嶺； 15 又臨到高臺和堅固城牆； 16 又臨到他施的船隻，並一切可愛的美物。 17 驕傲的必屈膝，狂妄的必降卑，在\underline{ 那日}，唯獨耶和華被尊崇。18 偶像必全然廢棄。 19 耶和華興起使地大震動的時候，人就進入石洞，進入土穴，躲避耶和華的驚嚇和他威嚴的榮光。 20 到那日，人必將為拜而造的金偶像、銀偶像拋給田鼠和蝙蝠。 21 到耶和華興起使地大震動的時候，人好進入磐石洞中和巖石穴裡，躲避耶和華的驚嚇和他威嚴的榮光。 22 你們休要倚靠世人，他鼻孔裡不過有氣息，他在一切事上可算什麼呢？


\subsubsection{以賽亞書 3：16-25}
 耶和華又說：「因為錫安的女子狂傲，行走挺項，賣弄眼目，俏步徐行，腳下玎璫， 17 所以主必使錫安的女子頭長禿瘡，耶和華又使她們赤露下體。」 18 到 \underline{那日}，主必除掉她們華美的腳釧、髮網、月牙圈、 19 耳環、手鐲、蒙臉的帕子、 20 華冠、足鏈、華帶、香盒、符囊、 21 戒指、鼻環、 22 吉服、外套、雲肩、荷包、 23 手鏡、細麻衣、裹頭巾、蒙身的帕子。 24 必有臭爛代替馨香，繩子代替腰帶，光禿代替美髮，麻衣繫腰代替華服，烙傷代替美容。 25 你的男丁必倒在刀下，你的勇士必死在陣上。 26 錫安的城門必悲傷哀號，她必荒涼坐在地上。

\subsubsection{以賽亞書 4：1-6}
1 在 \underline{那日}，七個女人必拉住一個男人，說：「我們吃自己的食物，穿自己的衣服，但求你許我們歸你名下，求你除掉我們的羞恥！」2 到 \underline{那日}，耶和華發生的苗必華美尊榮，地的出產必為以色列逃脫的人顯為榮華茂盛。 3 主以公義的靈和焚燒的靈，將錫安女子的汙穢洗去，又將耶路撒冷中殺人的血除淨。那時，剩在錫安留在耶路撒冷的，就是一切住耶路撒冷在生命冊上記名的，必稱為聖。 4  5 耶和華也必在錫安全山並各會眾以上，使白日有煙雲，黑夜有火焰的光，因為在全榮耀之上必有遮蔽。 6 必有亭子，白日可以得蔭避暑，也可以作為藏身之處，躲避狂風暴雨
\subsubsection{以賽亞書 5：24-30}
火苗怎樣吞滅碎秸，乾草怎樣落在火焰之中，照樣，他們的根必像朽物，他們的花必像灰塵飛騰。因為他們厭棄萬軍之耶和華的訓誨，藐視以色列聖者的言語。 25 所以耶和華的怒氣向他的百姓發作，他的手伸出攻擊他們，山嶺就震動，他們的屍首在街市上好像糞土。雖然如此，他的怒氣還未轉消，他的手仍伸不縮。26 他必豎立大旗招遠方的國民，發嘶聲叫他們從地極而來。看哪，他們必急速奔來！ 27 其中沒有疲倦的、絆跌的，沒有打盹的、睡覺的，腰帶並不放鬆，鞋帶也不折斷。 28 他們的箭快利，弓也上了弦；馬蹄算如堅石，車輪好像旋風。 29 他們要吼叫像母獅子，咆哮像少壯獅子；他們要咆哮抓食，坦然叼去，無人救回。 30 \underline{那日}，他們要向以色列人吼叫，像海浪砰訇。人若望地，只見黑暗艱難，光明在雲中變為昏暗。

\subsubsection{以賽亞書 6：9-13}
撒拉弗說：「你去告訴這百姓說：『你們聽是要聽見，卻不明白；看是要看見，卻不曉得。』 10 要使這百姓心蒙脂油，耳朵發沉，眼睛昏迷，恐怕眼睛看見，耳朵聽見，心裡明白，回轉過來，便得醫治。」 11 我就說：「主啊，這到幾時為止呢？」他說：「\underline{直到}城邑荒涼，無人居住，房屋空閒無人，地土極其荒涼。 12 並且耶和華將人遷到遠方，在這境內撇下的地土很多。 13 境內剩下的人，若還有十分之一，也必被吞滅。像栗樹、橡樹，雖被砍伐，樹不子卻仍存留，這聖潔的種類在國中也是如此。」

\subsubsection{以賽亞書 7：17-25}
17 耶和華必使亞述王\underline{攻擊你的日子}臨到你和你的百姓並你的父家，自從以法蓮離開猶大以來，\underline{未曾有這樣的日子}。18 「\underline{那時}，耶和華要發嘶聲，使埃及江河源頭的蒼蠅和亞述地的蜂子飛來。 19 都必飛來，落在荒涼的谷內，磐石的穴裡，和一切荊棘籬笆中，並一切的草場上。20 「那時，主必用大河外賃的剃頭刀，就是亞述王，剃去頭髮和腳上的毛，並要剃淨鬍鬚。21 「\underline{那時}，一個人要養活一隻母牛犢，兩隻母綿羊。 22 因為出的奶多，他就得吃奶油。在境內所剩的人，都要吃奶油與蜂蜜。23 「從前凡種一千棵葡萄樹，值銀一千舍客勒的地方，到那時，必長荊棘和蒺藜。 24 人上那裡去，必帶弓箭，因為遍地滿了荊棘和蒺藜。 25 所有用鋤刨挖的山地，你因怕荊棘和蒺藜，不敢上那裡去，只可成了放牛之處，為羊踐踏之地。」

\subsubsection{以賽亞書 8：5-22}
5 耶和華又曉諭我說： 6 「這百姓既厭棄西羅亞緩流的水，喜悅利汛和利瑪利的兒子， 7 因此主必使大河翻騰的水猛然沖來，就是亞述王和他所有的威勢。必漫過一切的水道，漲過兩岸； 8 必沖入猶大，漲溢氾濫，直到頸項。以馬內利啊！他展開翅膀，遍滿你的地。」9 列國的人民哪，任憑你們喧嚷，終必破壞！遠方的眾人哪，當側耳而聽！任憑你們束起腰來，終必破壞，你們束起腰來，終必破壞！ 10 任憑你們同謀，終歸無有，任憑你們言定，終不成立，因為神與我們同在。 11 耶和華以大能的手，指教我不可行這百姓所行的道，對我這樣說： 12 「這百姓說同謀背叛，你們不要說同謀背叛。他們所怕的，你們不要怕，也不要畏懼。 13 但要尊萬軍之耶和華為聖，以他為你們所當怕的，所當畏懼的。 14 他必作為聖所，卻向以色列兩家做絆腳的石頭、跌人的磐石，向耶路撒冷的居民作為圈套和網羅。 15 許多人必在其上絆腳跌倒，而且跌碎，並陷入網羅被纏住。」16 你要捲起律法書，在我門徒中間封住訓誨。 17 我要等候那掩面不顧雅各家的耶和華，我也要仰望他。 18 看哪，我與耶和華所給我的兒女，就是從住在錫安山萬軍之耶和華來的，在以色列中作為預兆和奇蹟.19 有人對你們說：「當求問那些交鬼的和行巫術的，就是聲音綿蠻、言語微細的。」你們便回答說：「百姓不當求問自己的神嗎？豈可為活人求問死人呢？ 20 人當以訓誨和法度為標準。」他們所說的若不與此相符，必不得見晨光。 21 他們必經過這地，受艱難，受飢餓，飢餓的時候心中焦躁，咒罵自己的君王和自己的神。 22 仰觀上天，俯察下地，不料，盡是艱難黑暗和幽暗的痛苦，他們必被趕入烏黑的黑暗中去。


\subsubsection{以賽亞書 9：11-21}
11 因此，耶和華要高舉利汛的敵人來攻擊以色列，並要激動以色列的仇敵， 12 東有亞蘭人，西有非利士人，他們張口要吞吃以色列。雖然如此，耶和華的怒氣還未轉消，他的手仍伸不縮。13 這百姓還沒有歸向擊打他們的主，也沒有尋求萬軍之耶和華。 14 因此，耶和華\underline{一日之間}，必從以色列中剪除頭與尾，棕枝與蘆葦。 15 長老和尊貴人就是頭，以謊言教人的先知就是尾。 16 因為引導這百姓的使他們走錯了路，被引導的都必敗亡。 17 所以主必不喜悅他們的少年人，也不憐恤他們的孤兒寡婦，因為各人是褻瀆的，是行惡的，並且各人的口都說愚妄的話。雖然如此，耶和華的怒氣還未轉消，他的手仍伸不縮。邪惡像火焚燒，燒滅荊棘和蒺藜，在稠密的樹林中著起來，就成為煙柱，旋轉上騰。 19 因萬軍之耶和華的烈怒，地都燒遍，百姓成為火柴，無人憐愛弟兄。 20 有人右邊搶奪，仍受飢餓；左邊吞吃，仍不飽足。各人吃自己膀臂上的肉。 21 瑪拿西吞吃以法蓮，以法蓮吞吃瑪拿西，又一同攻擊猶大。雖然如此，耶和華的怒氣還未轉消，他的手仍伸不縮

\subsubsection{以賽亞書 10：12-34}
 12 主在錫安山和耶路撒冷成就他一切工作的時候，主說：「我必罰亞述王自大的心和他高傲眼目的榮耀。 13 因為他說：『我所成就的事，是靠我手的能力和我的智慧，我本有聰明。我挪移列國的地界，搶奪他們所積蓄的財寶，並且我像勇士，使坐寶座的降為卑。 14 我的手夠到列國的財寶，好像人夠到鳥窩；我也得了全地，好像人拾起所棄的雀蛋。沒有動翅膀的，沒有張嘴的，也沒有鳴叫的。』」15 斧豈可向用斧砍木的自誇呢？鋸豈可向用鋸的自大呢？好比棍掄起那舉棍的，好比杖舉起那非木的人。 16 因此，主萬軍之耶和華必使亞述王的肥壯人變為瘦弱，在他的榮華之下必有火著起，如同焚燒一樣。 17 以色列的光必如火，他的聖者必如火焰，在一日之間，將亞述王的荊棘和蒺藜焚燒淨盡。 18 又將他樹林和肥田的榮耀全然燒盡，好像拿軍旗的昏過去一樣。 19 他林中剩下的樹必稀少，就是孩子也能寫其數。 到\underline{那日}，以色列所剩下的和雅各家所逃脫的，不再倚靠那擊打他們的，卻要誠實倚靠耶和華以色列的聖者。 21 所剩下的，就是雅各家所剩下的，必歸回全能的神。 22 以色列啊，你的百姓雖多如海沙，唯有剩下的歸回。原來滅絕的事已定，必有公義施行，如水漲溢。 23 因為主萬軍之耶和華在全地之中，必成就所定規的結局。24 所以主萬軍之耶和華如此說：「住錫安我的百姓啊，亞述王雖然用棍擊打你，又照埃及的樣子舉杖攻擊你，你卻不要怕他。 25 因為還有一點點時候，向你們發的憤恨就要完畢，我的怒氣要向他發作，使他滅亡。」 26 萬軍之耶和華要興起鞭來攻擊他，好像在俄立磐石那裡殺戮米甸人一樣；耶和華的杖要向海伸出，把杖舉起，像在埃及一樣。 27 到\underline{那日}，亞述王的重擔必離開你的肩頭，他的軛必離開你的頸項，那軛也必因肥壯的緣故撐斷。亞述王來到亞葉，經過米磯崙，在密抹安放輜重。 29 他們過了隘口，在迦巴住宿。拉瑪人戰兢，掃羅的基比亞人逃跑。 30 迦琳的居民 c哪，要高聲呼喊！萊煞人哪，須聽！哀哉，困苦的亞拿突啊！ 31 瑪得米那人躲避，基柄的居民逃遁。 32 當\underline{那日}，亞述王要在挪伯歇兵，向錫安女子的山，就是耶路撒冷的山，掄手攻他。33 看哪，主萬軍之耶和華以驚嚇削去樹枝，長高的必被砍下，高大的必被伐倒。 34 稠密的樹林，他要用鐵器砍下，黎巴嫩的樹木必被大能者伐倒。

\subsubsection{以賽亞書 11:1-5，10-16}
1 從耶西的本 a必發一條，從他根生的枝子必結果實。 2 耶和華的靈必住在他身上，就是使他有智慧和聰明的靈，謀略和能力的靈，知識和敬畏耶和華的靈。 3 他必以敬畏耶和華為樂，行審判不憑眼見，斷是非也不憑耳聞； 4 卻要以公義審判貧窮人，以正直判斷世上的謙卑人，以口中的杖擊打世界，以嘴裡的氣殺戮惡人。 5 公義必當他的腰帶，信實必當他脅下的帶子。到\underline{那日}，耶西的根立做萬民的大旗，外邦人必尋求他，他安息之所大有榮耀11 當\underline{那日}，主必二次伸手救回自己百姓中所餘剩的，就是在亞述、埃及、巴忒羅、古實、以攔、示拿、哈馬並眾海島所剩下的。 12 他必向列國豎立大旗，招回以色列被趕散的人，又從地的四方聚集分散的猶大人。 13 以法蓮的嫉妒就必消散，擾害猶大的必被剪除；以法蓮必不嫉妒猶大，猶大也不擾害以法蓮。 14 他們要向西飛，撲在非利士人的肩頭上，一同擄掠東方人，伸手按住以東和摩押，亞捫人也必順服他們。 15 耶和華必使埃及海汊枯乾，掄手用暴熱的風使大河分為七條，令人過去不致濕腳。 16 為主餘剩的百姓，就是從亞述剩下回來的，必有一條大道，如當日以色列從埃及地上來一樣。

\subsubsection{以賽亞書 13:2-18}
2 「應當在淨光的山豎立大旗，向群眾揚聲招手，使他們進入貴胄的門！ 3 我吩咐我所挑出來的人，我招呼我的勇士，就是那矜誇高傲之輩，為要成就我怒中所定的。」 4 山間有多人的聲音，好像是大國人民。有許多國的民聚集鬨嚷的聲音，這是萬軍之耶和華點齊軍隊，預備打仗。 5 他們從遠方來，從天邊來，就是耶和華並他惱恨的兵器，要毀滅這全地。6 你們要哀號，因為\underline{耶和華的日子}臨近了！這日來到，好像毀滅從全能者來到。 7 所以人手都必軟弱，人心都必消化。 8 他們必驚惶悲痛，愁苦必將他們抓住，他們疼痛好像產難的婦人一樣。彼此驚奇相看，臉如火焰。 9 \underline{耶和華的日子}臨到，必有殘忍、憤恨、烈怒，使這地荒涼，從其中除滅罪人。 10 天上的眾星群宿都不發光，日頭一出就變黑暗，月亮也不放光。 11 「我必因邪惡刑罰世界，因罪孽刑罰惡人，使驕傲人的狂妄止息，制伏強暴人的狂傲。 12 我必使人比精金還少，使人比俄斐純金更少。 13 我萬軍之耶和華在憤恨中發烈怒的日子，必使天震動，使地搖撼離其本位。 14 人必像被追趕的鹿，像無人收聚的羊，各歸回本族，各逃到本土。 15 凡被仇敵追上的必被刺死，凡被捉住的必被刀殺。 16 他們的嬰孩必在他們眼前摔碎，他們的房屋必被搶奪，他們的妻子必被玷汙。17 「我必激動瑪代人來攻擊他們，瑪代人不注重銀子，也不喜愛金子。 18 他們必用弓擊碎少年人，不憐憫婦人所生的，眼也不顧惜孩子。

\subsubsection{以賽亞書 17:4-9，12-14}
4 「到\underline{那日}，雅各的榮耀必致枵薄，他肥胖的身體必漸瘦弱。 5 就必像收割的人收斂禾稼，用手割取穗子，又像人在利乏音谷拾取遺落的穗子。 6 其間所剩下的不多，好像人打橄欖樹，在儘上的枝梢上只剩兩三個果子，在多果樹的旁枝上只剩四五個果子。」這是耶和華以色列的神說的。 7 當\underline{那日}，人必仰望造他們的主，眼目重看以色列的聖者。 8 他們必不仰望祭壇，就是自己手所築的，也不重看自己指頭所做的，無論是木偶是日像。 9 在\underline{那日}，他們的堅固城必像樹林中和山頂上所撇棄的地方，就是從前在以色列人面前被人撇棄的，這樣地就荒涼了....唉！多民鬨嚷，好像海浪砰訇；列邦奔騰，好像猛水滔滔。 13 列邦奔騰，好像多水滔滔，但神斥責他們，他們就遠遠逃避，又被追趕，如同山上的風前糠，又如暴風前的旋風土。 14 到\underline{晚上}有驚嚇，\underline{未到早晨}他們就沒有了。這是擄掠我們之人所得的份，是搶奪我們之人的報應。

\subsubsection{耶利米書 51:1-24}
耶和華如此說：「我必使毀滅的風颳起，攻擊巴比倫和住在立加米的人。 2 我要打發外邦人來到巴比倫，簸揚她，使她的地空虛，在她、\underline{遭禍的日子}他們要周圍攻擊她。 3 拉弓的要向拉弓的和貫甲挺身的射箭，不要憐惜她的少年人，要滅盡她的全軍。 4 他們必在迦勒底人之地被殺仆倒，在巴比倫的街上被刺透。」5 以色列和猶大雖然境內充滿違背以色列聖者的罪，卻沒有被他的神萬軍之耶和華丟棄。 6 你們要從巴比倫中逃奔，各救自己的性命！不要陷在她的罪孽中一同滅亡，因為這是耶和華報仇的時候，他必向巴比倫施行報應。 7 巴比倫素來是耶和華手中的金杯，使天下沉醉，萬國喝了她的酒就顛狂了。 8 巴比倫忽然傾覆毀壞，要為她哀號，為止她的疼痛，拿乳香或者可以治好。 9 我們想醫治巴比倫，她卻沒有治好。離開她吧！我們各人歸回本國，因為她受的審判通於上天，達到穹蒼。 10 耶和華已經彰顯我們的公義，來吧，我們可以在錫安報告耶和華我們神的作為。11 你們要磨尖了箭頭，抓住盾牌。耶和華定意攻擊巴比倫，將她毀滅，所以激動了瑪代君王的心，因這是耶和華報仇，就是為自己的殿報仇。 12 你們要豎立大旗攻擊巴比倫的城牆，要堅固瞭望臺，派定守望的設下埋伏。因為耶和華指著巴比倫居民所說的話、所定的意，他已經做成。 13 住在眾水之上多有財寶的啊，你的結局到了，你貪婪之量滿了！ 14 萬軍之耶和華指著自己起誓說：「我必使敵人充滿你，像螞蚱一樣，他們必呐喊攻擊你。」15 耶和華用能力創造大地，用智慧建立世界，用聰明鋪張穹蒼。 16 他一發聲，空中便有多水激動，他使雲霧從地極上騰。他造電隨雨而閃，從他府庫中帶出風來。 17 各人都成了畜類，毫無知識，各銀匠都因他的偶像羞愧。他所鑄的偶像本是虛假的，其中並無氣息， 18 都是虛無的，是迷惑人的工作，到追討的時候，必被除滅。 19 雅各的份不像這些，因他是造做萬有的主，以色列也是他產業的支派，萬軍之耶和華是他的名。20 「你是我爭戰的斧子和打仗的兵器。我要用你打碎列國，用你毀滅列邦， 21 用你打碎馬和騎馬的，用你打碎戰車和坐在其上的， 22 用你打碎男人和女人，用你打碎老年人和少年人，用你打碎壯丁和處女， 23 用你打碎牧人和他的群畜，用你打碎農夫和他一對牛，用你打碎省長和副省長。」 24 耶和華說：「我必在你們眼前，報復巴比倫人和迦勒底居民在錫安所行的諸惡。」25 耶和華說：「你這行毀滅的山哪，就是毀滅天下的山，我與你反對。我必向你伸手，將你從山巖滾下去，使你成為燒毀的山。 26 人必不從你那裡取石頭為房角石，也不取石頭為根基石，你必永遠荒涼。」這是耶和華說的。27 要在境內豎立大旗，在各國中吹角，使列國預備攻擊巴比倫，將亞拉臘、米尼、亞實基拿各國招來攻擊她。又派軍長來攻擊她，使馬匹上來如螞蚱。 28 使列國和瑪代君王，與省長和副省長，並他們所管全地之人都預備攻擊她。 29 地必震動而瘠苦，因耶和華向巴比倫所定的旨意成立了，使巴比倫之地荒涼，無人居住。 30 巴比倫的勇士止息爭戰，藏在堅壘之中，他們的勇力衰盡，好像婦女一樣。巴比倫的住處有火著起，門閂都折斷了。 31 跑報的要彼此相遇，送信的要互相迎接，報告巴比倫王說：「城的四方被攻取了， 32 渡口被占據了，葦塘被火燒了，兵丁也驚慌了。」33 萬軍之耶和華以色列的神如此說：「巴比倫城 a好像踹穀的禾場，再過片時，收割她的時候就到了。」 34 以色列人說：「巴比倫王尼布甲尼撒吞滅我，壓碎我，使我成為空虛的器皿。他像大魚將我吞下，用我的美物充滿他的肚腹，又將我趕出去。」 35 錫安的居民要說：「巴比倫以強暴待我，損害我的身體，願這罪歸給她！」耶路撒冷人要說：「願流我們血的罪歸到迦勒底的居民！」 36 所以，耶和華如此說：「我必為你申冤，為你報仇。我必使巴比倫的海枯竭，使她的泉源乾涸。 37 巴比倫必成為亂堆，為野狗的住處，令人驚駭、嗤笑，並且無人居住。 38 他們要像少壯獅子咆哮，像小獅子吼叫。 39 他們火熱的時候，我必為他們設擺酒席，使他們沉醉，好叫他們快樂，睡了長覺永不醒起。」這是耶和華說的。 40 「我必使他們像羊羔，像公綿羊和公山羊下到宰殺之地。41 「示沙克 b何竟被攻取！天下所稱讚的何竟被占據！巴比倫在列國中何竟變為荒場！ 42 海水漲起漫過巴比倫，她被許多海浪遮蓋。 43 她的城邑變為荒場、旱地、沙漠，無人居住、無人經過之地。 44 我必刑罰巴比倫的彼勒，使他吐出所吞的。萬民必不再流歸他那裡，巴比倫的城牆也必坍塌了。45 「我的民哪，你們要從其中出去，各人拯救自己，躲避耶和華的烈怒。 46 你們不要心驚膽怯，也不要因境內所聽見的風聲懼怕。因為這年有風聲傳來，那年也有風聲傳來，境內有強暴的事，官長攻擊官長。 47 \underline{日子}將到，我必刑罰巴比倫雕刻的偶像，她全地必然抱愧，她被殺的人必在其中仆倒。 48 \underline{那時}天地和其中所有的必因巴比倫歡呼，因為行毀滅的要從北方來到她那裡。」這是耶和華說的。 49 巴比倫怎樣使以色列被殺的人仆倒，照樣，她全地被殺的人也必在巴比倫仆倒。50 你們躲避刀劍的要快走，不要站住。要在遠方記念耶和華，心中追想耶路撒冷。 51 我們聽見辱罵就蒙羞，滿面慚愧，因為外邦人進入耶和華殿的聖所。 52 耶和華說：「\underline{日子}將到，我必刑罰巴比倫雕刻的偶像，通國受傷的人必唉哼。 53 巴比倫雖升到天上，雖使她堅固的高處更堅固，還有行毀滅的從我這裡到她那裡。」這是耶和華說的。54 有哀號的聲音從巴比倫出來，有大毀滅的響聲從迦勒底人之地發出。 55 因耶和華使巴比倫變為荒場，使其中的大聲滅絕。仇敵彷彿眾水，波浪砰訇，響聲已經發出。 56 這是行毀滅的臨到巴比倫，巴比倫的勇士被捉住，他們的弓折斷了，因為耶和華是施行報應的神，必定施行報應。 57 君王，名為萬軍之耶和華的說：「我必使巴比倫的首領、智慧人、省長、副省長和勇士都沉醉，使他們睡了長覺永不醒起。」 58 萬軍之耶和華如此說：「巴比倫寬闊的城牆必全然傾倒，她高大的城門必被火焚燒。眾民所勞碌的必致虛空，列國所勞碌的被火焚燒，他們都必困乏。」59 猶大王西底家在位第四年上巴比倫去的時候，瑪西雅的孫子、尼利亞的兒子西萊雅與王同去，西萊雅是王宮的大臣，先知耶利米有話吩咐他。 60 耶利米將一切要臨到巴比倫的災禍，就是論到巴比倫的一切話，寫在書上。 61 耶利米對西萊雅說：「你到了巴比倫，務要念這書上的話。 62 又說：『耶和華啊，你曾論到這地方說要剪除，甚至連人帶牲畜沒有在這裡居住的，必永遠荒涼。』




\subsubsection{以賽亞書 66:15-24}
「看哪，耶和華必在火中降臨，他的車輦像旋風，以烈怒施行報應，以火焰施行責罰。 16 因為耶和華在一切有血氣的人身上，必以火與刀施行審判，被耶和華所殺的必多。 17 那些分別為聖潔淨自己的，進入園內跟在其中一個人的後頭，吃豬肉和倉鼠並可憎之物，他們必一同滅絕。」這是耶和華說的。18 「我知道他們的行為和他們的意念。\underline{時候將到}，我必將萬民萬族聚來，看見我的榮耀。 19 我要顯神蹟在他們中間，逃脫的，我要差到列國去，就是到他施、普勒、拉弓的路德和土巴、雅完，並素來沒有聽見我名聲，沒有看見我榮耀遼遠的海島。他們必將我的榮耀傳揚在列國中。 20 他們必將你們的弟兄從列國中送回，使他們或騎馬，或坐車，坐轎，騎騾子，騎獨峰駝到我的聖山耶路撒冷，作為供物獻給耶和華，好像以色列人用潔淨的器皿盛供物奉到耶和華的殿中。」這是耶和華說的。 21 耶和華說：「我也必從他們中間取人為祭司，為利未人。」22 耶和華說：「我所要造的新天新地怎樣在我面前長存，你們的後裔和你們的名字也必照樣長存。 23 每逢月朔、安息日，凡有血氣的必來在我面前下拜。」這是耶和華說的。 24 「他們必出去觀看那些違背我人的屍首。因為他們的蟲是不死的，他們的火是不滅的。凡有血氣的，都必憎惡他們。」


\subsubsection{約珥書 2:1-21}
1 「你們要在錫安吹角，在我聖山吹出大聲！國中的居民都要發顫，因為\underline{耶和華的日子}將到，已經臨近。 2 \underline{那日}是黑暗、幽冥、密雲、烏黑的日子，好像晨光鋪滿山嶺。有一隊蝗蟲又大又強，從來沒有這樣的，以後直到萬代也必沒有。 3 牠們前面如火燒滅，後面如火焰燒盡。未到以前，地如伊甸園，過去以後，成了荒涼的曠野，沒有一樣能躲避牠們的。4 「牠們的形狀如馬，奔跑如馬兵。 5 在山頂蹦跳的響聲如車輛的響聲，又如火焰燒碎秸的響聲，好像強盛的民擺陣預備打仗。 6 牠們一來，眾民傷慟，臉都變色。 7 牠們如勇士奔跑，像戰士爬城，各都步行，不亂隊伍。 8 彼此並不擁擠，向前各行其路，直闖兵器，不偏左右。 9 牠們蹦上城，躥上牆，爬上房屋，進入窗戶如同盜賊。 10 牠們一來，地震天動，日月昏暗，星宿無光。 11 耶和華在他軍旅前發聲，他的隊伍甚大，成就他命的是強盛者。因為\underline{耶和華的日子}大而可畏，誰能當得起呢？2 「耶和華說：雖然如此，你們應當禁食、哭泣、悲哀，一心歸向我。 13 你們要撕裂心腸，不撕裂衣服，歸向耶和華你們的神，因為他有恩典，有憐憫，不輕易發怒，有豐盛的慈愛，並且後悔不降所說的災。 14 或者他轉意後悔，留下餘福，就是留下獻給耶和華你們神的素祭和奠祭，也未可知。15 「你們要在錫安吹角，分定禁食的日子，宣告嚴肅會。 16 聚集眾民，使會眾自潔，招聚老者，聚集孩童和吃奶的，使新郎出離洞房，新婦出離內室。 17 侍奉耶和華的祭司要在廊子和祭壇中間哭泣，說：『耶和華啊，求你顧惜你的百姓，不要使你的產業受羞辱，列邦管轄他們。為何容列國的人說：「他們的神在哪裡呢？」』「耶和華就為自己的地發熱心，憐恤他的百姓。 19 耶和華應允他的百姓說：我必賜給你們五穀、新酒和油，使你們飽足。我也不再使你們受列國的羞辱， 20 卻要使北方來的軍隊遠離你們，將他們趕到乾旱荒廢之地，前隊趕入東海，後隊趕入西海。因為他們所行的大惡 b，臭氣上升，腥味騰空。」21 地土啊，不要懼怕，要歡喜快樂！因為耶和華行了大事。 22 田野的走獸啊，不要懼怕！因為曠野的草發生，樹木結果，無花果樹、葡萄樹也都效力。 23 錫安的民哪，你們要快樂，為耶和華你們的神歡喜！因他賜給你們合宜的秋雨，為你們降下甘霖，就是秋雨、春雨，和先前一樣。 24 禾場必滿了麥子，酒榨與油榨必有新酒和油盈溢。 25 「我打發到你們中間的大軍隊，就是蝗蟲、蝻子、螞蚱、剪蟲，那些年所吃的，我要補還你們。 26 你們必多吃而得飽足，就讚美為你們行奇妙事之耶和華你們神的名。我的百姓必永遠不致羞愧。 27 你們必知道我是在以色列中間，又知道我是耶和華你們的神，在我以外並無別神。我的百姓必永遠不致羞愧。28 「以後，我要將我的靈澆灌凡有血氣的，你們的兒女要說預言，你們的老年人要做異夢，少年人要見異象。 29 在\underline{那些日子}，我要將我的靈澆灌我的僕人和使女。 30 在天上地下，我要顯出奇事，有血，有火，有煙柱。 31 日頭要變為黑暗，月亮要變為血，這都在\underline{耶和華大而可畏的日子}未到以前。32「\underline{到那時候}，凡求告耶和華名的，就必得救。因為照耶和華所說的，在錫安山耶路撒冷必有逃脫的人，在剩下的人中必有耶和華所召的。


\subsubsection{約珥書 3:1-21}
1 「到\underline{那日}，我使猶大和耶路撒冷被擄之人歸回的時候， 2 我要聚集萬民，帶他們下到約沙法谷，在那裡施行審判。因為他們將我的百姓，就是我的產業以色列，分散在列國中，又分取我的地土， 3 且為我的百姓拈鬮，將童子換妓女，賣童女買酒喝。 4 推羅、西頓和非利士四境的人哪，你們與我何干？你們要報復我嗎？若報復我，我必使報應速速歸到你們的頭上。 5 你們既然奪取我的金銀，又將我可愛的寶物帶入你們宮殿 a， 6 並將猶大人和耶路撒冷人賣給希臘人，使他們遠離自己的境界， 7 我必激動他們離開你們所賣到之地，又必使報應歸到你們的頭上。 8 我必將你們的兒女賣在猶大人的手中，他們必賣給遠方示巴國的人。這是耶和華說的。9 「當在萬民中宣告說：要預備打仗，激動勇士，使一切戰士上前來！ 10 要將犁頭打成刀劍，將鐮刀打成戈矛。軟弱的要說：『我有勇力！』」 11 四圍的列國啊，你們要速速地來，一同聚集！耶和華啊，求你使你的大能者降臨！ 12 「萬民都當興起，上到約沙法谷，因為我必坐在那裡，審判四圍的列國。 13 開鐮吧！因為莊稼熟了。踐踏吧！因為酒榨滿了，酒池盈溢，他們的罪惡甚大。14 「許多許多的人在斷定谷，因為耶和華的日子臨近斷定谷。 15 日月昏暗，星宿無光。 16 耶和華必從錫安吼叫，從耶路撒冷發聲，天地就震動。耶和華卻要做他百姓的避難所，做以色列人的保障。 17 你們就知道我是耶和華你們的神，且又住在錫安我的聖山。那時，耶路撒冷必成為聖，外邦人不再從其中經過。18 「到\underline{那日}，大山要滴甜酒，小山要流奶子，猶大溪河都有水流，必有泉源從耶和華的殿中流出來，滋潤什亭谷。 19 埃及必然荒涼，以東變為淒涼的曠野，都因向猶大人所行的強暴，又因在本地流無辜人的血。「但猶大必存到永遠，耶路撒冷必存到萬代。 21 我未曾報復流血的罪，現在我要報復，因為耶和華住在錫安。」

\subsubsection{撒迦利亞書 14:1-21}

1 \underline{耶和華的日子}臨近，你的財物必被搶掠，在你中間分散。 2 「因為我必聚集萬國與耶路撒冷爭戰，城必被攻取，房屋被搶奪，婦女被玷汙，城中的民一半被擄去，剩下的民仍在城中，不致剪除。」 3 那時耶和華必出去與那些國爭戰，好像從前爭戰一樣。4 \underline{那日}，他的腳必站在耶路撒冷前面朝東的橄欖山上。這山必從中間分裂，自東至西成為極大的谷，山的一半向北挪移，一半向南挪移。 5 「你們要從我山的谷中逃跑，因為山谷必延到亞薩。你們逃跑，必如猶大王烏西雅年間的人逃避大地震一樣。」耶和華我的神必降臨，有一切聖者同來。 6 \underline{那日}必沒有光，三光必退縮。 7 \underline{那日}必是耶和華所知道的，不是白晝，也不是黑夜，到了晚上才有光明。 8 \underline{那日}必有活水從耶路撒冷出來，一半往東海流，一半往西海流，冬夏都是如此。
9 耶和華必做全地的王，那日耶和華必為獨一無二的，他的名也是獨一無二的。 10 全地，從迦巴直到耶路撒冷南方的臨門，要變為亞拉巴。耶路撒冷必仍居高位，就是從便雅憫門到第一門之處，又到角門，並從哈楠業樓直到王的酒榨。 11 人必住在其中，不再有咒詛，耶路撒冷人必安然居住。12 耶和華用災殃攻擊那與耶路撒冷爭戰的列國人，必是這樣：他們兩腳站立的時候，肉必消沒，眼在眶中乾癟，舌在口中潰爛。 13 \underline{那日}，耶和華必使他們大大擾亂，他們各人彼此揪住，舉手攻擊。 14 猶大也必在耶路撒冷爭戰。那時四圍各國的財物，就是許多金銀、衣服，必被收聚。 15 那臨到馬匹、騾子、駱駝、驢和營中一切牲畜的災殃是與那災殃一般。16 所有來攻擊耶路撒冷列國中剩下的人，必年年上來敬拜大君王萬軍之耶和華，並守住棚節。 17 地上萬族中，凡不上耶路撒冷敬拜大君王萬軍之耶和華的，必無雨降在他們的地上。 18 埃及族若不上來，雨也不降在他們的地上。凡不上來守住棚節的列國人，耶和華也必用這災攻擊他們。 19 這就是埃及的刑罰和那不上來守住棚節之列國的刑罰。 20 當\underline{那日}，馬的鈴鐺上必有「歸耶和華為聖的」這句話。耶和華殿內的鍋必如祭壇前的碗一樣。 21 凡耶路撒冷和猶大的鍋都必歸萬軍之耶和華為聖，凡獻祭的都必來取這鍋，煮肉在其中。當\underline{那日}，在萬軍之耶和華的殿中必不再有迦南人。


\subsubsection{瑪拉基書 3:1-4，16-18}
萬軍之耶和華說：「我要差遣我的使者在我前面預備道路。你們所尋求的主必忽然進入他的殿，立約的使者，就是你們所仰慕的，快要來到。「他來的日子，誰能當得起呢？他顯現的時候，誰能立得住呢？因為他如煉金之人的火，如漂布之人的鹼。 3 他必坐下如煉淨銀子的，必潔淨利未人，熬煉他們像金銀一樣，他們就憑公義獻供物給耶和華。 4 \underline{那時}猶大和耶路撒冷所獻的供物必蒙耶和華悅納，彷彿古時之日，上古之年。」16  那時，敬畏耶和華的彼此談論，耶和華側耳而聽，且有紀念冊在他面前，記錄那敬畏耶和華思念他名的人。 17 萬軍之耶和華說：「在我\underline{所定的日子}他們必屬我，特特歸我，我必憐恤他們，如同人憐恤服侍自己的兒子。 18 那時你們必歸回，將善人和惡人，侍奉神的和不侍奉神的，分別出來。

\subsubsection{瑪拉基書 4:1-5}
1 萬軍之耶和華說：「\underline{那日}臨近，勢如燒著的火爐，凡狂傲的和行惡的必如碎秸，在\underline{那日}必被燒盡，根本、枝條一無存留。 2 但向你們敬畏我名的人，必有公義的日頭出現，其光線有醫治之能。你們必出來，跳躍如圈裡的肥犢。 3 你們必踐踏惡人，在我\underline{所定的日子}，他們必如灰塵在你們腳掌之下。這是萬軍之耶和華說的。4 「你們當記念我僕人摩西的律法，就是我在何烈山為以色列眾人所吩咐他的律例、典章。「看哪，\underline{耶和華大而可畏之日}未到以前，我必差遣先知以利亞到你們那裡去。 6 他必使父親的心轉向兒女，兒女的心轉向父親，免得我來咒詛遍地。」


\subsubsection{使徒行傳 2:14-22}
14 彼得和十一個使徒站起，高聲說：「猶太人和一切住在耶路撒冷的人哪，這件事你們當知道，也當側耳聽我的話。 15 你們想這些人是醉了，其實不是醉了，因為時候剛到巳初。 16 這正是先知約珥所說的： 17 『神說：在\underline{末後的日子}，我要將我的靈澆灌凡有血氣的，你們的兒女要說預言，你們的少年人要見異象，老年人要做異夢。 18 在\underline{那些日子}，我要將我的靈澆灌我的僕人和使女，他們就要說預言。 19 在天上我要顯出奇事，在地下我要顯出神蹟，有血，有火，有煙霧。 20 日頭要變為黑暗，月亮要變為血，這都在\underline{主大而明顯的日子}未到以前。 21 到那時候，凡求告主名的，就必得救。』 22 以色列人哪，請聽我的話：神藉著拿撒勒人耶穌在你們中間施行異能、奇事、神蹟，將他證明出來，


\subsubsection{帖撒羅尼迦後書 1:3-10}
3 弟兄們，我們該為你們常常感謝神，這本是合宜的，因你們的信心格外增長，並且你們眾人彼此相愛的心也都充足。 4 甚至我們在神的各教會裡為你們誇口，都因你們在所受的一切逼迫患難中，仍舊存忍耐和信心。 5 這正是神公義判斷的明證，叫你們可算配得神的國，你們就是為這國受苦。 6 神既是公義的，就必將患難報應那加患難給你們的人， 7 也必使你們這受患難的人與我們同得平安。那時，主耶穌同他有能力的天使從天上在火焰中顯現， 8 要報應那不認識神和那不聽從我主耶穌福音的人。 9 他們要受刑罰，就是永遠沉淪，離開主的面和他權能的榮光。 10 這正是主降臨，要在他聖徒的身上得榮耀，又在一切信的人身上顯為稀奇的\underline{那日子}。我們對你們作的見證，你們也信了。


\subsubsection{彼得後書 3:3-13}
 3 第一要緊的，該知道在末世必有好譏誚的人隨從自己的私慾出來譏誚說： 4 「主要降臨的應許在哪裡呢？因為從列祖睡了以來，萬物與起初創造的時候仍是一樣。」 5 他們故意忘記：從太古憑神的命有了天，並從水而出、藉水而成的地； 6 故此，當時的世界被水淹沒就消滅了。 7 但現在的天地還是憑著那命存留，直留到不敬虔之人受審判遭沉淪的日子，用火焚燒。8 親愛的弟兄啊，有一件事你們不可忘記，就是主看一日如千年，千年如一日。 9 主所應許的尚未成就，有人以為他是耽延，其實不是耽延，乃是寬容你們，不願有一人沉淪，乃願人人都悔改。 10 但主的日子要像賊來到一樣。那日，天必大有響聲廢去，有形質的都要被烈火銷化，地和其上的物都要燒盡了。 11 這一切既然都要如此銷化，你們為人該當怎樣聖潔，怎樣敬虔， 12 切切仰望神的日子來到！在\underline{那日}，天被火燒就銷化了，有形質的都要被烈火熔化。 13 但我們照他的應許，盼望新天新地，有義居在其中。

\subsection{大水}

\subsubsection{詩篇 18：1-19}
1 耶和華我的力量啊，我愛你！
2 耶和華是我的巖石，我的山寨，我的救主，我的神，我的磐石，我所投靠的。他是我的盾牌，是拯救我的角，是我的高臺。
3 我要求告當讚美的耶和華，這樣我必從仇敵手中被救出來。
4 曾有死亡的繩索纏繞我，匪類的\underline{急流}使我驚懼，
5 陰間的繩索纏繞我，死亡的網羅臨到我。
6 我在急難中求告耶和華，向我的神呼求。他從殿中聽了我的聲音，我在他面前的呼求入了他的耳中。
7 那時因他發怒，地就搖撼戰抖，山的根基也震動搖撼。
8 從他鼻孔冒煙上騰，從他口中發火焚燒，連炭也著了。
9 他又使天下垂，親自降臨，有黑雲在他腳下。
他又使天下垂，親自降臨，有黑雲在他腳下。
10 他坐著基路伯飛行，他藉著風的翅膀快飛。
11 他以黑暗為藏身之處，以水的黑暗、天空的厚雲為他四圍的行宮。
12 因他面前的光輝，他的厚雲行過，便有冰雹火炭。
13 耶和華也在天上打雷，至高者發出聲音，便有冰雹火炭。
14 他射出箭來，使仇敵四散；多多發出閃電，使他們擾亂。
15 耶和華啊，你的斥責一發，你鼻孔的氣一出，海底就出現，大地的根基也顯露。
16 他從高天伸手抓住我，把我從\underline{大水}中拉上來。
17 他救我脫離我的勁敵和那些恨我的人，因為他們比我強盛。
18 我遭遇災難的日子，他們來攻擊我，但耶和華是我的倚靠。
19 他又領我到寬闊之處；他救拔我，因他喜悅我。


\subsubsection{詩篇 69:14}
求你搭救我出離淤泥，不叫我陷在其中；求你使我脫離那些恨我的人，使我出離\underline{深水}。

\subsubsection{詩篇 124：1-7}
1 以色列人要說：「若不是耶和華幫助我們，
2 若不是耶和華幫助我們，當人起來攻擊我們，
3 向我們發怒的時候，就把我們活活地吞了。
4 那時\underline{波濤}必漫過我們，\underline{河水}必淹沒我們，
5 \underline{狂傲的水}必淹沒我們。」
6 耶和華是應當稱頌的！他沒有把我們當野食交給他們吞吃。
7 我們好像雀鳥從捕鳥人的網羅裡逃脫，網羅破裂，我們逃脫了。
8 我們得幫助，是在乎倚靠造天地之耶和華的名。

\subsubsection{詩篇 144：7}
7 求你從上伸手救拔我，救我出離\underline{大水}，救我脫離外邦人的手。

\subsubsection{以賽亞書 17：12-14}
唉！多民鬨嚷，好像\underline{海浪砰訇}；列邦奔騰，好像\underline{猛水滔滔}。 13 列邦奔騰，好像\underline{多水滔滔}，但神斥責他們，他們就遠遠逃避，又被追趕，如同山上的風前糠，又如暴風前的旋風土。 14 到晚上有驚嚇，未到早晨他們就沒有了。這是擄掠我們之人所得的份，是搶奪我們之人的報應。

\subsubsection{耶利米書 51：54-58}
54 有哀號的聲音從巴比倫出來，有大毀滅的響聲從迦勒底人之地發出。 55 因耶和華使巴比倫變為荒場，使其中的大聲滅絕。仇敵彷彿\underline{眾水，波浪砰訇}，響聲已經發出。 56 這是行毀滅的臨到巴比倫，巴比倫的勇士被捉住，他們的弓折斷了，因為耶和華是施行報應的神，必定施行報應。 57 君王，名為萬軍之耶和華的說：「我必使巴比倫的首領、智慧人、省長、副省長和勇士都沉醉，使他們睡了長覺永不醒起。」 58 萬軍之耶和華如此說：「巴比倫寬闊的城牆必全然傾倒，她高大的城門必被火焚燒。眾民所勞碌的必致虛空，列國所勞碌的被火焚燒，他們都必困乏。」


\subsection{巴比倫}

\subsubsection{耶利米書 51：54-58}
54 有哀號的聲音從\underline{巴比倫}出來，有大毀滅的響聲從迦勒底人之地發出。 55 因耶和華使\underline{巴比倫}變為荒場，使其中的大聲滅絕。仇敵彷彿眾水，波浪砰訇，響聲已經發出。 56 這是行毀滅的臨到\underline{巴比倫}，巴比倫的勇士被捉住，他們的弓折斷了，因為耶和華是施行報應的神，必定施行報應。 57 君王，名為萬軍之耶和華的說：「我必使\underline{巴比倫}的首領、智慧人、省長、副省長和勇士都沉醉，使他們睡了長覺永不醒起。」 58 萬軍之耶和華如此說：「\underline{巴比倫}寬闊的城牆必全然傾倒，她高大的城門必被火焚燒。眾民所勞碌的必致虛空，列國所勞碌的被火焚燒，他們都必困乏。」


\subsection{末期}
\subsubsection{但以理書 12：5-13}
5 我但以理觀看，見另有兩個人站立，一個在河這邊，一個在河那邊。 6 有一個問那站在河水以上穿細麻衣的說：「這奇異的事到幾時才應驗呢？」 7 我聽見那站在河水以上穿細麻衣的，向天舉起左右手，指著活到永遠的主起誓說：「要到一載，二載，半載，打破聖民權力的時候，這一切事就都應驗了。」 8 我聽見這話，卻不明白，就說：「我主啊，這些事的\underline{結局אַחֲרִ֖ית}是怎樣呢？」 9 他說：「但以理啊，你只管去，因為這話已經隱藏封閉，直到\underline{קֵץ末時}。 10 必有許多人使自己清淨潔白，且被熬煉，但惡人仍必行惡。一切惡人都不明白，唯獨智慧人能明白。 11 從除掉常獻的燔祭，並設立那行毀壞可憎之物的時候，必有一千二百九十日。 12 等到一千三百三十五日的，那人便為有福！ 13 你且去等候\underline{結局}，因為你必安歇。到了\underline{末期קֵץ}，你必起來，享受你的福分。」

\subsubsection{馬太福音 10：5-42}
5 耶穌差這十二個人去，吩咐他們說：「外邦人的路你們不要走，撒馬利亞人的城你們不要進， 6 寧可往以色列家迷失的羊那裡去。 7 隨走隨傳，說『天國近了』。 8 醫治病人，叫死人復活，叫長大痲瘋的潔淨，把鬼趕出去。你們白白地得來，也要白白地捨去。 9 腰袋裡不要帶金銀銅錢， 10 行路不要帶口袋，不要帶兩件褂子，也不要帶鞋和拐杖，因為工人得飲食是應當的。 11 你們無論進哪一城、哪一村，要打聽那裡誰是好人，就住在他家，直住到走的時候。 12 進他家裡去，要請他的安。 13 那家若配得平安，你們所求的平安就必臨到那家；若不配得，你們所求的平安仍歸你們。 14 凡不接待你們、不聽你們話的人，你們離開那家或是那城的時候，就把腳上的塵土跺下去。 15 我實在告訴你們：當審判的日子，所多瑪和蛾摩拉所受的比那城還容易受呢！21 「弟兄要把弟兄，父親要把兒子，送到死地；兒女要與父母為敵，害死他們。 22 並且你們要為我的名被眾人恨惡，唯有忍耐到\underline{底τέλος}的，必然得救。 23 有人在這城裡逼迫你們，就逃到那城裡去。我實在告訴你們：以色列的城邑你們還沒有走遍，人子就到了。24 「學生不能高過先生，僕人不能高過主人。 25 學生和先生一樣，僕人和主人一樣，也就罷了。人既罵家主是別西卜 a，何況他的家人呢？ 26 所以不要怕他們，因為掩蓋的事沒有不露出來的，隱藏的事沒有不被人知道的。 27 我在暗中告訴你們的，你們要在明處說出來；你們耳中所聽的，要在房上宣揚出來。 28 那殺身體不能殺靈魂的，不要怕他們；唯有能把身體和靈魂都滅在地獄裡的，正要怕他。 29 兩個麻雀不是賣一分銀子嗎？若是你們的父不許，一個也不能掉在地上。 30 就是你們的頭髮，也都被數過了。 31 所以，不要懼怕，你們比許多麻雀還貴重。 32 凡在人面前認我的，我在我天上的父面前也必認他； 33 凡在人面前不認我的，我在我天上的父面前也必不認他。34 「你們不要想我來是叫地上太平，我來並不是叫地上太平，乃是叫地上動刀兵。 35 因為我來是叫『人與父親生疏，女兒與母親生疏，媳婦與婆婆生疏； 36 人的仇敵就是自己家裡的人』。 37 愛父母過於愛我的，不配做我的門徒；愛兒女過於愛我的，不配做我的門徒； 38 不背著他的十字架跟從我的，也不配做我的門徒。 39 得著生命的，將要失喪生命；為我失喪生命的，將要得著生命.40 「人接待你們，就是接待我；接待我，就是接待那差我來的。 41 人因為先知的名接待先知，必得先知所得的賞賜；人因為義人的名接待義人，必得義人所得的賞賜。 42 無論何人，因為門徒的名，只把一杯涼水給這小子裡的一個喝，我實在告訴你們：這人不能不得賞賜。」



\subsubsection{馬太福音 24：4-14}
 4 耶穌回答說：「你們要謹慎，免得有人迷惑你們。 5 因為將來有好些人冒我的名來，說：『我是基督』，並且要迷惑許多人。 6 你們也要聽見打仗和打仗的風聲，總不要驚慌，因為這些事是必須有的，只是\underline{末期 τέλος}還沒有到。 7 民要攻打民，國要攻打國，多處必有饑荒、地震， 8 這都是災難的起頭。 9 那時，人要把你們陷在患難裡，也要殺害你們；你們又要為我的名被萬民恨惡。10 「那時，必有許多人跌倒，也要彼此陷害、彼此恨惡， 11 且有好些假先知起來，迷惑多人。 12 只因不法的事增多，許多人的愛心才漸漸冷淡了。 13 唯有忍耐到\underline{底τέλος}的，必然得救。 14 這天國的福音要傳遍天下，對萬民作見證，然後\underline{末期 τέλος}才來到。

\subsubsection{馬可福音 13：3-13}
3 耶穌在橄欖山上對聖殿而坐，彼得、雅各、約翰和安得烈暗暗地問他說： 4 「請告訴我們，什麼時候有這些事呢？這一切事將成的時候有什麼預兆呢？」 5 耶穌說：「你們要謹慎，免得有人迷惑你們。 6 將來有好些人冒我的名來，說『我是基督』，並且要迷惑許多人。 7 你們聽見打仗和打仗的風聲，不要驚慌，這些事是必須有的，只是\underline{末期τέλος}還沒有到。 8 民要攻打民，國要攻打國，多處必有地震、饑荒，這都是災難的起頭。9 「但你們要謹慎，因為人要把你們交給公會，並且你們在會堂裡要受鞭打，又為我的緣故站在諸侯與君王面前，對他們作見證。 10 然而，福音必須先傳給萬民。 11 人把你們拉去交官的時候，不要預先思慮說什麼，到那時候，賜給你們什麼話，你們就說什麼；因為說話的不是你們，乃是聖靈。 12 弟兄要把弟兄，父親要把兒子，送到死地；兒女要起來與父母為敵，害死他們。 13 並且你們要為我的名被眾人恨惡，唯有忍耐到\underline{底τέλος}的，必然得救。

\subsubsection{路加福音 21：6-9}
7 他們問他說：「夫子，什麼時候有這事呢？這事將到的時候有什麼預兆呢？」 8 耶穌說：「你們要謹慎，不要受迷惑。因為將來有好些人冒我的名來，說『我是基督』，又說『時候近了』。你們不要跟從他們！ 9 你們聽見打仗和擾亂的事，不要驚惶，因為這些事必須先有，只是\underline{末期τέλος}不能立時就到。」
\subsubsection{路加福音 21：20-28}
20 「你們看見耶路撒冷被兵圍困，就可知道它成荒場的日子近了。 21 那時，在猶太的，應當逃到山上；在城裡的，應當出來；在鄉下的，不要進城。 22 因為這是報應的日子，使經上所寫的都得應驗。 23 當那些日子，懷孕的和奶孩子的有禍了！因為將有大災難降在這地方，也有震怒臨到這百姓。 24 他們要倒在刀下，又被擄到各國去。耶路撒冷要被外邦人踐踏，直到外邦人的日期滿了。 25 日、月、星辰要顯出異兆，地上的邦國也有困苦，因海中波浪的響聲就慌慌不定。 26 天勢都要震動，人想起那將要臨到世界的事，就都嚇得魂不附體。 27 那時，他們要看見人子有能力、有大榮耀駕雲降臨。 28 一有這些事，你們就當挺身昂首，因為你們\underline{得贖的日子}近了。」

\subsubsection{哥林多前書 1：4-9}
4 我常為你們感謝我的神，因神在基督耶穌裡所賜給你們的恩惠， 5 又因你們在他裡面凡事富足，口才、知識都全備。 6 正如我為基督作的見證在你們心裡得以堅固， 7 以致你們在恩賜上沒有一樣不及人的，等候我們的主耶穌基督顯現。 8 他也必堅固你們到\underline{底τέλους}，叫你們在我們主耶穌基督的日子無可責備。 9 神是信實的，你們原是被他所召，好與他兒子我們的主耶穌基督一同得份。


\subsubsection{哥林多前書 10：4-13}
1 弟兄們，我不願意你們不曉得：我們的祖宗從前都在雲下，都從海中經過， 2 都在雲裡、海裡受洗歸了摩西， 3 並且都吃了一樣的靈食， 4 也都喝了一樣的靈水；所喝的是出於隨著他們的靈磐石，那磐石就是基督。 5 但他們中間多半是神不喜歡的人，所以在曠野倒斃。 6 這些事都是我們的鑒戒，叫我們不要貪戀惡事，像他們那樣貪戀的。 7 也不要拜偶像，像他們有人拜的，如經上所記：「百姓坐下吃喝，起來玩耍。」 8 我們也不要行姦淫，像他們有人行的，一天就倒斃了二萬三千人。 9 也不要試探主，像他們有人試探的，就被蛇所滅。 10 你們也不要發怨言，像他們有發怨言的，就被滅命的所滅。 11 他們遭遇這些事都要作為鑒戒，並且寫在經上正是警戒我們這\underline{末世τέλη}的人。 12 所以，自己以為站得穩的須要謹慎，免得跌倒。 13 你們所遇見的試探，無非是人所能受的。神是信實的，必不叫你們受試探過於所能受的，在受試探的時候，總要給你們開一條出路，叫你們能忍受得住。

\subsubsection{哥林多前書 15：20-28}
20 但基督已經從死裡復活，成為睡了之人初熟的果子。 21 死既是因一人而來，死人復活也是因一人而來。 22 在亞當裡眾人都死了，照樣，在基督裡眾人也都要復活。 23 但各人是按著自己的次序復活：初熟的果子是基督，以後在他來的時候，是那些屬基督的。 24 再後\underline{末期τέλος}到了，那時基督既將一切執政的、掌權的、有能的都毀滅了，就把國交於父神。 25 因為基督必要做王，等神把一切仇敵都放在他的腳下。 26 儘末了所毀滅的仇敵就是死， 27 因為經上說：「神叫萬物都服在他的腳下。」既說萬物都服了他，明顯那叫萬物服他的不在其內了。 28 萬物既服了他，那時子也要自己服那叫萬物服他的，叫神在萬物之上、為萬物之主。

\subsubsection{希伯來書 6：1-8}
1 所以，我們應當離開基督道理的開端，竭力進到完全的地步，不必再立根基，就如那懊悔死行、信靠神、 2 各樣洗禮、按手之禮、死人復活，以及永遠審判各等教訓。 3 神若許我們，我們必如此行。 4 論到那些已經蒙了光照，嘗過天恩的滋味，又於聖靈有份， 5 並嘗過神善道的滋味，覺悟來世權能的人， 6 若是離棄道理，就不能叫他們重新懊悔了。因為他們把神的兒子重釘十字架，明明地羞辱他。 7 就如一塊田地，吃過屢次下的雨水，生長菜蔬，合乎耕種的人用，就從神得福； 8 若長荊棘和蒺藜，必被廢棄，近於咒詛，\underline{結局 τέλους}就是焚燒。


\subsubsection{彼得前書 4：7-19}
7 萬物的\underline{結局τέλος}近了，所以你們要謹慎自守，警醒禱告。 8 最要緊的是彼此切實相愛，因為愛能遮掩許多的罪。 9 你們要互相款待，不發怨言。 10 各人要照所得的恩賜彼此服侍，做神百般恩賜的好管家。 11 若有講道的，要按著神的聖言講；若有服侍人的，要按著神所賜的力量服侍，叫神在凡事上因耶穌基督得榮耀——原來榮耀、權能都是他的，直到永永遠遠！阿們。12 親愛的弟兄啊，有火煉的試驗臨到你們，不要以為奇怪，似乎是遭遇非常的事； 13 倒要歡喜，因為你們是與基督一同受苦，使你們在他榮耀顯現的時候，也可以歡喜快樂。 14 你們若為基督的名受辱罵，便是有福的，因為神榮耀的靈常住在你們身上。 15 你們中間卻不可有人因為殺人、偷竊、作惡、好管閒事而受苦。 16 若為做基督徒受苦，卻不要羞恥，倒要因這名歸榮耀給神。 17 因為時候到了，審判要從神的家起首。若是先從我們起首，那不信從神福音的人將有何等的\underline{結局τέλος}呢？ 18 「若是義人僅僅得救，那不虔敬和犯罪的人將有何地可站呢？」 19 所以，那照神旨意受苦的人要一心為善，將自己靈魂交於那信實的造化之主。

\subsubsection{啟示錄 2：24-29}

24 『至於你們推雅推喇其餘的人，就是一切不從那教訓、不曉得他們素常所說撒旦深奧之理的人，我告訴你們：我不將別的擔子放在你們身上； 25 但你們已經有的，總要持守，直等到\underline{（the end）τέλους}我來。 26 那得勝又遵守我命令到底的，我要賜給他權柄制伏列國， 27 他必用鐵杖轄管 b他們，將他們如同窯戶的瓦器打得粉碎，像我從我父領受的權柄一樣。 28 我又要把晨星賜給他。 29 聖靈向眾教會所說的話，凡有耳的，就應當聽！』

\subsubsection{啟示錄 21：6-8}
他又對我說：「都成了！我是阿拉法，我是俄梅戛；我是初，我是\underline{終 τέλος}。我要將生命泉的水白白賜給那口渴的人喝。 7 得勝的，必承受這些為業：我要做他的神，他要做我的兒子。 8 唯有膽怯的、不信的、可憎的、殺人的、淫亂的、行邪術的、拜偶像的和一切說謊話的，他們的份就在燒著硫磺的火湖裡，這是第二次的死。」

\subsubsection{啟示錄 22：12-15}
「看哪，我必快來！賞罰在我，要照各人所行的報應他。 13 我是阿拉法，我是俄梅戛；我是首先的，我是\underline{末後τέλος}的；我是初，我是終。 14 那些洗淨自己衣服的有福了！可得權柄能到生命樹那裡，也能從門進城。 15 城外有那些犬類、行邪術的、淫亂的、殺人的、拜偶像的，並一切喜好說謊言、編造虛謊的。


\subsection{被提726 $\grave{\alpha}$$\rho$$\pi
$$\acute{\alpha}$$\zeta$$\omega$}

\subsubsection{馬太福音 11：11-19 $\grave{\alpha}$$\rho$$\pi$$\acute{\alpha}$$\zeta$$o$$\upsilon$$\sigma$$\iota$$\nu$}
「我實在告訴你們：凡婦人所生的，沒有一個興起來大過施洗約翰的；然而天國裡最小的比他還大。 12 從施洗約翰的時候到如今，天國是努力進入的，努力的人就\underline{得著$\grave{\alpha}$$\rho$$\pi$$\acute{\alpha}$$\zeta$$o$$\upsilon$$\sigma$$\iota$$\nu$}了。 13 因為眾先知和律法說預言到約翰為止； 14 你們若肯領受，這人就是那應當來的以利亞。 15 有耳可聽的，就應當聽！ 16 我可用什麼比這世代呢？好像孩童坐在街市上招呼同伴說： 17 『我們向你們吹笛，你們不跳舞！我們向你們舉哀，你們不捶胸！』 18 約翰來了，也不吃也不喝，人就說他是被鬼附著的。 19 人子來了，也吃也喝，人又說他是貪食好酒的人，是稅吏和罪人的朋友。但智慧之子總以智慧為是。」

\subsubsection{馬太福音 12：25-30 $\grave{\alpha}$$\rho$$\pi$$\acute{\alpha}$$\sigma$$\alpha$$\iota$}
25 耶穌知道他們的意念，就對他們說：「凡一國自相紛爭，就成為荒場；一城一家自相紛爭，必站立不住。 26 若撒旦趕逐撒旦，就是自相紛爭，他的國怎能站得住呢？ 27 我若靠著別西卜趕鬼，你們的子弟趕鬼又靠著誰呢？這樣，他們就要斷定你們的是非。 28 我若靠著神的靈趕鬼，這就是神的國臨到你們了。 29 人怎能進壯士家裡\underline{搶奪$\grave{\alpha}$$\rho$$\pi$$\acute{\alpha}$$\sigma$$\alpha$$\iota$ἁρπάσαι}他的家具呢？除非先捆住那壯士，才可以搶奪他的家財。 30 不與我相合的，就是敵我的；不同我收聚的，就是分散的。

\subsubsection{使徒行傳 8：36-40 $\acute{\eta}$$\rho$$\pi$$\alpha$$\sigma$$\epsilon$$\nu$}
36 二人正往前走，到了有水的地方，太監說：「看哪，這裡有水！我受洗有什麼妨礙呢？」 38 於是吩咐車站住，腓利和太監二人同下水裡去，腓利就給他施洗。 39 從水裡上來，主的靈把腓利\underline{提 $\acute{\eta}$$\rho$$\pi$$\alpha$$\sigma$$\epsilon$$\nu$ ἥρπασεν}了去。太監也不再見他了，就歡歡喜喜地走路。 40 後來有人在亞鎖都遇見腓利；他走遍那地方，在各城宣傳福音，直到愷撒利亞。

\subsubsection{哥林多後書 12：1-4 $\grave{\alpha}$$\rho$$\pi$$\alpha$$\gamma$$\acute{\epsilon}$$\nu$$\tau$$\alpha$, $\grave{\eta}$$\rho$$\pi$$\acute{\alpha}$$\gamma$$\eta$}
1 我自誇固然無益，但我是不得已的。如今我要說到主的顯現和啟示。 2 我認得一個在基督裡的人，他前十四年被\underline{提$\grave{\alpha}$$\rho$$\pi$$\alpha$$\gamma$$\acute{\epsilon}$$\nu$$\tau$$\alpha$ἁρπαγέντα}到第三層天上去——或在身內，我不知道；或在身外，我也不知道；只有神知道。 3 我認得這人——或在身內、或在身外，我都不知道，只有神知道—— 4 他被\underline{提$\grave{\eta}$$\rho$$\pi$$\acute{\alpha}$$\gamma$$\eta$ἡρπάγη}到樂園裡，聽見隱祕的言語，是人不可說的。

\subsubsection{帖撒羅尼迦前書 4：13-18 $\grave{\alpha}$$\rho$$\pi$$\alpha$$\gamma$$\eta$$\sigma$$\acute{o}$$\mu$$\epsilon$$\theta$$\alpha$}
論到睡了的人，我們不願意弟兄們不知道，恐怕你們憂傷，像那些沒有指望的人一樣。 14 我們若信耶穌死而復活了，那已經在耶穌裡睡了的人，神也必將他們與耶穌一同帶來。 15 我們現在照主的話告訴你們一件事：我們這活著還存留到主降臨的人，斷不能在那已經睡了的人之先。 16 因為主必親自從天降臨，有呼叫的聲音和天使長的聲音，又有神的號吹響，那在基督裡死了的人必先復活， 17 以後我們這活著還存留的人必和他們一同被\underline{提 $\grave{\alpha}$$\rho$$\pi$$\alpha$$\gamma$$\eta$$\sigma$$\acute{o}$$\mu$$\epsilon$$\theta$$\alpha$ἁρπαγησόμεθα}到雲裡，在空中與主相遇；這樣，我們就要和主永遠同在。 18 所以，你們當用這些話彼此勸慰。

\subsubsection{啟示錄 12：3-6 $\grave{\eta}$$\rho$$\pi$$\acute{\alpha}$$\sigma$$\theta$$\eta$}
3 天上又現出異象來：有一條大紅龍，七頭十角，七頭上戴著七個冠冕， 4 牠的尾巴拖拉著天上星辰的三分之一，摔在地上。龍就站在那將要生產的婦人面前，等她生產之後，要吞吃她的孩子。 5 婦人生了一個男孩子，是將來要用鐵杖轄管萬國的；她的孩子被\underline{提 $\grave{\eta}$$\rho$$\pi$$\acute{\alpha}$$\sigma$$\theta$$\eta$}到神寶座那裡去了。 6 婦人就逃到曠野，在那裡有神給她預備的地方，使她被養活一千二百六十天。

\subsection{奥秘}
\subsubsection{但以理書 2:22, 29-30, 46-48}
\begin{itemize} 
\itemsep-.25em
\item 2:22 他顯明\underline{深奧隱祕的事}，知道暗中所有的，光明也與他同居。
\item 2:29 王啊，你在床上想到後來的事，那顯明\underline{奧祕事}的主把將來必有的事指示你。至於那\underline{奧祕}的事顯明給我，並非因我的智慧勝過一切活人，乃為使王知道夢的講解和心裡的思念。
\item 46 當時尼布甲尼撒王俯伏在地，向但以理下拜，並且吩咐人給他奉上供物和香品。 47王對但以理說：「你既能顯明這\underline{奧祕}的事，你們的神誠然是萬神之神，萬王之主，又是顯明\underline{奧祕}事的。」 48於是王高抬但以理，賞賜他許多上等禮物，派他管理巴比倫全省，又立他為總理，掌管巴比倫的一切哲士
\end{itemize}

\subsubsection{阿摩司書 3:1-7}
1 「以色列人哪，你們全家是我從埃及地領上來的，當聽耶和華攻擊你們的話： 2 在地上萬族中我只認識你們，因此，我必追討你們的一切罪孽。 3 二人若不同心，豈能同行呢？ 4 獅子若非抓食，豈能在林中咆哮呢？少壯獅子若無所得，豈能從洞中發聲呢？ 5 若沒有機檻，雀鳥豈能陷在網羅裡呢？網羅若無所得，豈能從地上翻起呢？ 6 城中若吹角，百姓豈不驚恐呢？災禍若臨到一城，豈非耶和華所降的嗎？ 7 主耶和華若不將\underline{奧祕}指示他的僕人眾先知，就一無所行


\subsubsection{馬太福音 13:10-12, 13:35}
\begin{itemize} 
\itemsep-.25em
\item 10門徒進前來，問耶穌說：「對眾人講話為什麼用比喻呢？」 11耶穌回答說：「因為\underline{天國的奧祕}只叫你們知道，不叫他們知道。 12凡有的，還要加給他，叫他有餘；凡沒有的，連他所有的也要奪去。
\item 13:35 這是要應驗先知的話說：「我要開口用比喻，把創世以來所\underline{隱藏的事}發明出來。」
\end{itemize} 

\subsubsection{路加福音 8:10}
他說：「\underline{神國的奧祕}只叫你們知道，至於別人，就用比喻，叫他們『看也看不見，聽也聽不明』。

\subsubsection{羅馬書 2:16, 11:25, 16:25}
\begin{itemize} 
\itemsep-.25em
\item 就在神藉耶穌基督審判人\underline{隱祕事}的日子，照著我的福音所言。
\item 11:25 弟兄們，我不願意你們不知道這\underline{奧祕}，恐怕你們自以為聰明，就是：以色列人有幾分是硬心的，等到外邦人的數目添滿了，
\item 16:25 唯有神能照我所傳的福音和所講的耶穌基督，並照\underline{永古隱藏不言的奧祕}，堅固你們的心。
\end{itemize} 

\subsubsection{哥林多前書 2:1-13, 4:1-5, 14:1-3}
\begin{itemize} 
\itemsep-.25em
\item 1 弟兄們，從前我到你們那裡去，並沒有用高言大智對你們宣傳\underline{神的奧祕}。 2 因為我曾定了主意，在你們中間不知道別的，只知道耶穌基督並他釘十字架。 3 我在你們那裡，又軟弱，又懼怕，又甚戰兢。 4 我說的話、講的道，不是用智慧委婉的言語，乃是用聖靈和大能的明證， 5 叫你們的信不在乎人的智慧，只在乎神的大能。6 然而，在完全的人中，我們也講智慧，但不是這世上的智慧，也不是這世上有權有位將要敗亡之人的智慧。 7 我們講的，乃是\underline{從前所隱藏、神奧祕的智慧}，就是神在萬世以前預定使我們得榮耀的。 8 這智慧，世上有權有位的人沒有一個知道的，他們若知道，就不把榮耀的主釘在十字架上了。 9 如經上所記：「神為愛他的人所預備的，是眼睛未曾看見，耳朵未曾聽見，人心也未曾想到的。」 10 只有神藉著聖靈向我們顯明了。因為聖靈參透萬事，就是\underline{神深奧的事}也參透了。 11 除了在人裡頭的靈，誰知道人的事？像這樣，除了神的靈，也沒有人知道神的事。 12 我們所領受的，並不是世上的靈，乃是從神來的靈，叫我們能知道神開恩賜給我們的事。 13 並且我們講說這些事，不是用人智慧所指教的言語，乃是用聖靈所指教的言語，將屬靈的話解釋屬靈的事 。
\item 4:1 人應當以我們為基督的執事，為\underline{神奧祕事}的管家。 2 所求於管家的，是要他有忠心。 3 我被你們論斷，或被別人論斷，我都以為極小的事，連我自己也不論斷自己。 4 我雖不覺得自己有錯，卻也不能因此得以稱義，但判斷我的乃是主。 5 所以，時候未到，什麼都不要論斷，只等主來，他要照出暗中的隱情，顯明人心的意念。那時，各人要從神那裡得著稱讚。
\item 14:1 你們要追求愛，也要切慕屬靈的恩賜，其中更要羨慕的是做先知講道。 2那說方言的原不是對人說，乃是對神說，因為沒有人聽出來，然而他在心靈裡卻是講說各樣的\underline{奧祕}。 3但做先知講道的是對人說，要造就、安慰、勸勉人。
\end{itemize} 

\subsubsection{以弗所書 1:3-12, 3:3-4, 14:1-3}
\begin{itemize} 
\itemsep-.25em
\item 1:3 他在基督裡曾賜給我們天上各樣屬靈的福氣， 4 就如神從創立世界以前在基督裡揀選了我們，使我們在他面前成為聖潔，無有瑕疵； 5 又因愛我們，就按著自己意旨所喜悅的，預定我們藉著耶穌基督得兒子的名分， 6 使他榮耀的恩典得著稱讚。這恩典是他在愛子裡所賜給我們的。 7 我們藉這愛子的血得蒙救贖，過犯得以赦免，乃是照他豐富的恩典。 8 這恩典是神用諸般智慧聰明，充充足足賞給我們的， 9 都是照他自己所預定的美意，叫我們知道他\underline{旨意的奧祕}， 10 要照所安排的，在日期滿足的時候，使天上、地上一切所有的，都在基督裡面同歸於一。 11 我們也在他裡面得基業，這原是那位隨己意行做萬事的，照著他旨意所預定的， 12 叫他的榮耀從我們這首先在基督裡有盼望的人，可以得著稱讚。
\item 3:1 因此，我保羅——為你們外邦人做了基督耶穌被囚的，替你們祈禱 。 2 諒必你們曾聽見神賜恩給我，將關切你們的職分託付我， 3 用啟示使我知道\underline{福音的奧祕} ，正如我以前略略寫過的。你們念了，就能曉得我深知\underline{基督的奧祕} 。5 這\underline{奧祕}在以前的世代沒有叫人知道，像如今藉著聖靈啟示他的聖使徒和先知一樣。 6 這\underline{奧祕}就是外邦人在基督耶穌裡，藉著福音，得以同為後嗣，同為一體，同蒙應許。 7 我做了這福音的執事，是照神的恩賜，這恩賜是照他運行的大能賜給我的。 8我本來比眾聖徒中最小的還小，然而他還賜我這恩典，叫我把基督那測不透的豐富傳給外邦人， 9又使眾人都明白，這歷代以來隱藏在創造萬物之神\underline{裡的奧祕} 是如何安排的， 10為要藉著教會，使天上執政的、掌權的現在得知神百般的智慧。 11 這是照神從萬世以前在我們主基督耶穌裡所定的旨意。 12 我們因信耶穌，就在他裡面放膽無懼，篤信不疑地來到神面前。 13 所以，我求你們不要因我為你們所受的患難喪膽，這原是你們的榮耀。
\item  5:22 你們做妻子的，當順服自己的丈夫，如同順服主； 23 因為丈夫是妻子的頭，如同基督是教會的頭，他又是教會全體的救主。 24 教會怎樣順服基督，妻子也要怎樣凡事順服丈夫。 25 你們做丈夫的，要愛你們的妻子，正如基督愛教會，為教會捨己， 26 要用水藉著道把教會洗淨，成為聖潔， 27 可以獻給自己，做個榮耀的教會，毫無玷汙、皺紋等類的病，乃是聖潔沒有瑕疵的。 28 丈夫也當照樣愛妻子，如同愛自己的身子，愛妻子便是愛自己了。 29 從來沒有人恨惡自己的身子，總是保養顧惜，正像基督待教會一樣， 30 因我們是他身上的肢體 。 31 為這個緣故，人要離開父母，與妻子聯合，二人成為一體。 32這是\underline{極大的奧祕} ，但我是指著基督和教會說的。 33然而你們各人都當愛妻子，如同愛自己一樣；妻子也當敬重她的丈夫。
\item 6:18靠著聖靈，隨時多方禱告祈求，並要在此警醒不倦，為眾聖徒祈求； 19也為我祈求，使我得著口才，能以放膽開口講明\underline{福音的奧祕}—— 20我為這\underline{福音的奧祕} 做了戴鎖鏈的使者——並使我照著當盡的本分放膽講論。
\end{itemize} 

\subsubsection{歌羅西書 1:25-27}
25我照神為你們所賜我的職分做了教會的執事，要把神的道理傳得全備。 26這道理就是歷世歷代所隱藏的奧祕，但如今向他的聖徒顯明了。 27神願意叫他們知道，這奧祕在外邦人中有何等豐盛的榮耀，就是基督在你們心裡成了有榮耀的盼望。
歌羅西書 2:2
要叫他們的心得安慰，因愛心互相聯絡，以致豐豐足足在悟性中有充足的信心，使他們真知神的奧祕，就是基督，
歌羅西書 3:3
因為你們已經死了，你們的生命與基督一同藏在神裡面。
歌羅西書 4:3
也要為我們禱告，求神給我們開傳道的門，能以講基督的奧祕——我為此被捆鎖——

提摩太前書 3:9
8做執事的也是如此，必須端莊，不一口兩舌，不好喝酒，不貪不義之財； 9要存清潔的良心，固守真道的奧祕。 10這等人也要先受試驗，若沒有可責之處，然後叫他們做執事

啟示錄 10:7
但在第七位天使吹號發聲的時候，神的奧祕就成全了，正如神所傳給他僕人眾先知的佳音。」
啟示錄 11:15
第七位天使吹號，天上就有大聲音說：「世上的國成了我主和主基督的國，他要做王，直到永永遠遠。」

啟示錄 17:7
7天使對我說：「你為什麼稀奇呢？我要將這女人和馱著她的那七頭十角獸的奧祕告訴你。 8你所看見的獸，先前有，如今沒有，將要從無底坑裡上來，又要歸於沉淪。凡住在地上，名字從創世以來沒有記在生命冊上的，見先前有、如今沒有、以後再有的獸，就必稀奇



\subsection{生命册和羔羊生命册}

\subsubsection{出埃及記 32:30-35}
30 到了第二天，摩西對百姓說：「你們犯了大罪，我如今要上耶和華那裡去，或者可以為你們贖罪。」 31 摩西回到耶和華那裡，說：「唉！這百姓犯了大罪，為自己做了金像。 32 倘或你肯赦免他們的罪——不然，求你從\underline{你所寫的冊上}塗抹我的名。」 33 耶和華對摩西說：「誰得罪我，我就從\underline{我的冊上}塗抹誰的名。 34 現在你去領這百姓，往我所告訴你的地方去，我的使者必在你前面引路，只是到我追討的日子，我必追討他們的罪。」 35 耶和華殺百姓的緣故是因他們同亞倫做了牛犢。

\subsubsection{詩篇 69:21-28，56:8}
\begin{itemize} 
\itemsep-.35em 
\item 21 他們拿苦膽給我當食物，我渴了，他們拿醋給我喝。
22 願他們的筵席在他們面前變為網羅，在他們平安的時候變為機檻。
23 願他們的眼睛昏矇，不得看見，願你使他們的腰常常戰抖。
24 求你將你的惱恨倒在他們身上，叫你的烈怒追上他們。
25 願他們的住處變為荒場，願他們的帳篷無人居住。
26 因為你所擊打的，他們就逼迫；你所擊傷的，他們戲說他的愁苦。
27 願你在他們的罪上加罪，不容他們在你面前稱義。
28 願他們從\underline{生命冊}上被塗抹，不得\underline{記錄在義人之中}。
\item 我幾次流離，你都記數，求你把我眼淚裝在你的皮袋裡。這不都記在\underline{你冊子}上嗎
\end{itemize}

\subsubsection{以賽亞書 4:2-6}
2 到那日，耶和華發生的苗必華美尊榮，地的出產必為以色列逃脫的人顯為榮華茂盛。 3 主以公義的靈和焚燒的靈，將錫安女子的汙穢洗去，又將耶路撒冷中殺人的血除淨。那時，剩在錫安留在耶路撒冷的，就是一切住耶路撒冷在\underline{生命冊}上記名的，必稱為聖。 4  5 耶和華也必在錫安全山並各會眾以上，使白日有煙雲，黑夜有火焰的光，因為在全榮耀之上必有遮蔽。 6 必有亭子，白日可以得蔭避暑，也可以作為藏身之處，躲避狂風暴雨。


\subsubsection{耶利米書 17:13}
耶和華以色列的盼望啊，凡離棄你的，必致蒙羞。耶和華說：「離開我的，他們的\underline{名字必寫在土裡}，因為他們離棄我這活水的泉源。」

\subsubsection{以西結書 13:9}
我的手必攻擊那見虛假異象，用謊詐占卜的先知。他們必不列在我百姓的會中，不錄在\underline{以色列家的冊上}，也不進入以色列地，你們就知道我是主耶和華。』



\subsubsection{但以理書 12:1-4}
1 「那時，保佑你本國之民的天使長米迦勒必站起來，並且有大艱難，從有國以來直到此時沒有這樣的。你本國的民中，凡\underline{名錄在冊上}的，必得拯救。 2 睡在塵埃中的必有多人復醒，其中有得永生的，有受羞辱永遠被憎惡的。 3 智慧人必發光，如同天上的光；那使多人歸義的，必發光如星，直到永永遠遠。 4 但以理啊，你要隱藏這話，封閉這書，直到末時。必有多人來往奔跑 ，知識就必增長。」

\subsubsection{瑪拉基書 3:16-18}
那時，敬畏耶和華的彼此談論，耶和華側耳而聽，且有\underline{紀念冊}在他面前，記錄那敬畏耶和華思念他名的人。 17 萬軍之耶和華說：「在我所定的日子他們必屬我，特特歸我，我必憐恤他們，如同人憐恤服侍自己的兒子。 18 那時你們必歸回，將善人和惡人，侍奉神的和不侍奉神的，分別出來。」

\subsubsection{路加福音 10：17-20}
 那七十個人歡歡喜喜地回來，說：「主啊，因你的名，就是鬼也服了我們！」 18 耶穌對他們說：「我曾看見撒旦從天上墜落，像閃電一樣。 19 我已經給你們權柄可以踐踏蛇和蠍子，又勝過仇敵一切的能力，斷沒有什麼能害你們。 20 然而，不要因鬼服了你們就歡喜，要因你們的名\underline{記錄在天上}歡喜。」

 \subsubsection{腓立比書 4：2-3}
 2 我勸尤阿迪婭和循都基要在主裡同心。 3 我也求你這真實同負一軛的，幫助這兩個女人；因為她們在福音上曾與我一同勞苦，還有革利免並其餘和我一同做工的，他們的名字都在\underline{生命冊}上。
 
  \subsubsection{希伯來書 12：18-23}
18 你們原不是來到那能摸的山，此山有火焰、密雲、黑暗、暴風、 19 角聲與說話的聲音。那些聽見這聲音的，都求不要再向他們說話， 20 因為他們當不起所命他們的話說：「靠近這山的，即便是走獸，也要用石頭打死。」 21 所見的極其可怕，甚至摩西說：「我甚是恐懼戰兢。」 22 你們乃是來到錫安山，永生神的城邑，就是天上的耶路撒冷；那裡有千萬的天使， 23 有\underline{名錄在天上}諸長子之會所共聚的總會，有審判眾人的神和被成全之義人的靈魂， 24 並新約的中保耶穌以及所灑的血。這血所說的比亞伯的血所說的更美
 
 
 \subsubsection{生命冊 - 啟示錄 3:4-6, 13:8，17:8, 20:11-15} 
 \begin{itemize} 
\itemsep-.25em
\item 4然而在撒狄，你還有幾名是未曾汙穢自己衣服的，他們要穿白衣與我同行，因為他們是配得過的。 5凡得勝的，必這樣穿白衣，我也必不從\underline{生命冊}上塗抹他的名，且要在我父面前和我父眾使者面前，認他的名。 6聖靈向眾教會所說的話，凡有耳的，就應當聽！』	
\item 13:8 凡住在地上，名字從創世以來沒有記在被殺之\underline{羔羊生命冊}上的人，都要拜牠。
\item 你所看見的獸，先前有，如今沒有，將要從無底坑裡上來，又要歸於沉淪。凡住在地上，名字從創世以來沒有記在\underline{生命冊}上的，見先前有、如今沒有、以後再有的獸，就必稀奇。
\item 11 我又看見一個白色的大寶座與坐在上面的，從他面前天地都逃避，再無可見之處了。 12 我又看見死了的人，無論大小，都站在寶座前，\underline{案卷}展開了；並且另有\underline{一卷}展開，就是\underline{生命冊}。死了的人都憑著\underline{這些案卷}所記載的，照他們所行的受審判。 13 於是海交出其中的死人，死亡和陰間也交出其中的死人，他們都照各人所行的受審判。 14 死亡和陰間也被扔在火湖裡。這火湖就是第二次的死。 15 若有人名字沒記在\underline{生命冊}上，他就被扔在火湖裡。
\end{itemize} 

\subsubsection{羔羊生命册 - 啟示錄 13:7-10, 21:21-27}
\begin{itemize} 
\itemsep-.25em
\item 7 又任憑牠與聖徒爭戰，並且得勝；也把權柄賜給牠，制伏各族、各民、各方、各國。 8 凡住在地上，名字從創世以來沒有記在被殺之\underline{羔羊生命冊}上的人，都要拜牠。 9 凡有耳的，就應當聽！ 10 擄掠人的必被擄掠，用刀殺人的必被刀殺。聖徒的忍耐和信心就是在此。
\item 21 十二個門是十二顆珍珠，每門是一顆珍珠。城內的街道是精金，好像明透的玻璃。 22 我未見城內有殿，因主神全能者和羔羊為城的殿。 23 那城內又不用日月光照，因有神的榮耀光照，又有羔羊為城的燈。 24 列國要在城的光裡行走，地上的君王必將自己的榮耀歸於那城。 25 城門白晝總不關閉，在那裡原沒有黑夜。 26 人必將列國的榮耀、尊貴歸於那城。 27 凡不潔淨的，並那行可憎與虛謊之事的，總不得進那城，只有名字寫在\underline{羔羊生命冊}上的才得進去。
\end{itemize} 
